\chapter{Boolean Analysis — Decision Tree Fourier}
%%% generated by the blueprint pipeline from TCSlib.BooleanAnalysis.LMN.DecisionTreeFourier
\section{Overview}

This module develops the Fourier spectrum of a function computed by a decision tree
(O'Donnell Proposition 3.16). A recursive coefficient function computes the Fourier
expansion of the $\pm 1$-encoded tree function by structural induction, using the branch
identity $f = (f_{\mathrm{lo}} + f_{\mathrm{hi}})/2 + \chi_i (f_{\mathrm{lo}} -
f_{\mathrm{hi}})/2$ together with $\chi_i \chi_S = \chi_{S \triangle \{i\}}$. From this
recursion the module reads off the degree bound, the spectral $1$-norm bound, granularity
of the coefficients, and sparsity of the Fourier support, and restates the degree bound in
terms of the minimum decision-tree depth $\mathrm{dtDepth}$ as needed downstream in the
LMN pipeline.

%%%%%%%%%%%%%%%%%%%%%%%%%%%%%%%%%%%%%%%%%%%%%%%%%%%%%%%%%%%%%%%%%%%%%%%%%%%%%%%
\section{Declarations}

\begin{definition}[Size of a decision tree]
\lean{DecisionTree.size}
\label{DecisionTree.size}
\leanok
\statementsource{boolean-ch03-spectral-learning-restrictions}{pdf-20e15a3d3d5e-p006-b001}
\uses{DecisionTree}
The size of a decision tree, defined as its number of leaves: a leaf has size $1$, and a
branch has size equal to the sum of the sizes of its two subtrees.
\end{definition}

\begin{definition}[Sign-encoded tree function]
\lean{DecisionTree.signEval}
\label{DecisionTree.signEval}
\leanok
\statementsource{boolean-ch03-spectral-learning-restrictions}{pdf-20e15a3d3d5e-p005-b002}
\uses{BooleanAnalysis.BooleanFunc, BooleanAnalysis.boolToSign, DecisionTree, DecisionTree.eval}
The $\pm 1$-valued function computed by a decision tree $T$, namely
$x \mapsto \mathrm{boolToSign}(T.\mathrm{eval}\,x)$, where $\mathrm{false} \mapsto 1$ and
$\mathrm{true} \mapsto -1$.
\end{definition}

\begin{definition}[Recursive Fourier coefficients of a tree]
\lean{DecisionTree.coeffs}
\label{DecisionTree.coeffs}
\leanok
\uses{BooleanAnalysis.boolToSign, DecisionTree}
The coefficient function $\mathrm{coeffs}\,T : \mathcal{P}(\mathrm{Fin}\,n) \to \bbr$
defined by structural recursion on $T$. A leaf labelled $b$ assigns
$\mathrm{boolToSign}(b)$ to the empty frequency and $0$ to all others; a branch on
variable $i$ with subtrees $\mathrm{lo}, \mathrm{hi}$ assigns to $S$ the value
\[
\frac{\mathrm{coeffs}\,\mathrm{lo}\,S + \mathrm{coeffs}\,\mathrm{hi}\,S}{2}
+ \frac{\mathrm{coeffs}\,\mathrm{lo}\,(S \triangle \{i\})
      - \mathrm{coeffs}\,\mathrm{hi}\,(S \triangle \{i\})}{2}.
\]
\end{definition}

\begin{lemma}[Symmetric difference with a singleton is an involution]
\lean{DecisionTree.symmDiff_singleton_invol}
\label{DecisionTree.symmDiff_singleton_invol}
\leanok
\difficulty{1}
Fix $n \in \bbn$ and an index $i \in \{0, 1, \dots, n-1\}$. The map on subsets of $\{0,
1, \dots, n-1\}$ that sends each set $S$ to its symmetric difference $S \triangle \{i\}$
with the singleton $\{i\}$ is an involution: applying it twice returns every set to
itself, so $\bigl(S \triangle \{i\}\bigr) \triangle \{i\} = S$ for all $S$.
\end{lemma}

\begin{lemma}[Invariance of a subset-sum under symmetric difference with a point]
\lean{DecisionTree.sum_symmDiff_reindex}
\label{DecisionTree.sum_symmDiff_reindex}
\leanok
\uses{DecisionTree.symmDiff_singleton_invol}
\difficulty{1}
Fix $n \in \bbn$, let $i$ be an element of $\{0, 1, \dots, n-1\}$, and let $g$ assign a
real number to each subset of $\{0, 1, \dots, n-1\}$. Then summing $g$ over all subsets
is unaffected by first toggling the membership of $i$; that is,
\[
\sum_{S} g\big(S \mathbin{\triangle} \{i\}\big) = \sum_{S} g(S),
\]
where both sums range over all subsets $S \subseteq \{0, 1, \dots, n-1\}$ and $S
\mathbin{\triangle} \{i\}$ denotes the symmetric difference of $S$ with the singleton
$\{i\}$.
\end{lemma}

\begin{lemma}[Symmetric difference by a singleton toggles a Walsh character]
\lean{DecisionTree.chiS_symmDiff_singleton}
\label{DecisionTree.chiS_symmDiff_singleton}
\leanok
\uses{BooleanAnalysis.BoolCube, BooleanAnalysis.boolToSign, BooleanAnalysis.chiS, BooleanAnalysis.chiS_mul_chiS, BooleanAnalysis.chiS_singleton}
\difficulty{1}
Let $S\subseteq[n]$, let $i\in[n]$, and let $x\in\{0,1\}^n$. Then the Walsh--Fourier
character indexed by the symmetric difference $S\mathbin{\triangle}\{i\}$ satisfies
\[
  \chi_{S\,\triangle\,\{i\}}(x) \;=\; \chi_S(x)\,\sigma(x_i),
\]
where $\sigma$ is the sign encoding $\sigma(x_i)=(-1)^{x_i}$.
\end{lemma}

\begin{lemma}[Cardinality under singleton symmetric difference]
\lean{DecisionTree.card_symmDiff_singleton}
\label{DecisionTree.card_symmDiff_singleton}
\leanok
\difficulty{4}
Let $S$ be a finite subset of $\{0,1,\dots,n-1\}$ and let $i$ be one of these indices.
Then the symmetric difference $S \triangle \{i\}$ satisfies
\[
\abs{S} - 1 \le \abs{S \triangle \{i\}},
\]
where $\abs{S} - 1$ denotes truncated subtraction on the natural numbers (so it is $0$
when $S$ is empty).
\end{lemma}

\begin{lemma}[Walsh expansion of a decision tree's sign function]
\lean{DecisionTree.signEval_eq_sum_coeffs}
\label{DecisionTree.signEval_eq_sum_coeffs}
\leanok
\uses{BooleanAnalysis.BoolCube, BooleanAnalysis.boolToSign, BooleanAnalysis.chiS, BooleanAnalysis.chiS_empty, DecisionTree, DecisionTree.chiS_symmDiff_singleton, DecisionTree.coeffs, DecisionTree.eval, DecisionTree.signEval, DecisionTree.sum_symmDiff_reindex}
\difficulty{6}
Let $T$ be a decision tree on $n$ Boolean variables, computing the $\pm 1$-valued
sign-encoded function $x \mapsto \sigma(T(x))$, where $\sigma$ is the sign encoding
sending $\mathrm{false}$ to $1$ and $\mathrm{true}$ to $-1$. For every point $x \in
\{0,1\}^n$,
\[
  \sigma(T(x)) \;=\; \sum_{S \subseteq [n]} c_T(S)\,\chi_S(x),
\]
where the sum ranges over all subsets $S$ of $[n]$, $c_T(S)$ denotes the recursively
defined coefficient of $T$ at the frequency $S$, and $\chi_S$ is the Walsh character of
$S$.
\end{lemma}

\begin{lemma}[Fourier coefficients vanish above the tree depth]
\lean{DecisionTree.coeffs_eq_zero_of_depth_lt}
\statementsource{boolean-ch03-spectral-learning-restrictions}{pdf-20e15a3d3d5e-p006-b003}
\label{DecisionTree.coeffs_eq_zero_of_depth_lt}
\leanok
\uses{DecisionTree, DecisionTree.card_symmDiff_singleton, DecisionTree.coeffs, DecisionTree.depth}
\difficulty{4}
Let $T$ be a decision tree on $n$ Boolean variables, and let $S \subseteq
\{0,1,\dots,n-1\}$ be a set of variables whose cardinality $\abs{S}$ exceeds the depth
of $T$. Then the coefficient assigned to the frequency $S$ vanishes,
$\mathrm{coeffs}\,T\,S = 0$.
\end{lemma}

\begin{lemma}[Coefficient one-norm bounded by tree size]
\lean{DecisionTree.sum_abs_coeffs_le}
\statementsource{boolean-ch03-spectral-learning-restrictions}{pdf-20e15a3d3d5e-p006-b003}
\label{DecisionTree.sum_abs_coeffs_le}
\leanok
\uses{BooleanAnalysis.boolToSign, DecisionTree, DecisionTree.coeffs, DecisionTree.size, DecisionTree.sum_symmDiff_reindex}
\difficulty{5}
Let $T$ be a decision tree on $n$ Boolean variables, and for each subset $S \subseteq
\{0,\dots,n-1\}$ let $\mathrm{coeffs}\,T\,(S)$ denote its recursively defined Fourier
coefficient at frequency $S$. Then the total absolute mass of these coefficients is
bounded by the number of leaves of $T$:
\[
\sum_{S} \abs{\mathrm{coeffs}\,T\,(S)} \;\le\; \mathrm{size}(T).
\]
\end{lemma}

\begin{lemma}[Dyadic granularity of decision-tree Fourier coefficients]
\lean{DecisionTree.coeffs_mul_two_pow_int}
\statementsource{boolean-ch03-spectral-learning-restrictions}{pdf-20e15a3d3d5e-p006-b003}
\label{DecisionTree.coeffs_mul_two_pow_int}
\leanok
\uses{BooleanAnalysis.boolToSign, DecisionTree, DecisionTree.coeffs, DecisionTree.depth}
\difficulty{5}
Let $T$ be a decision tree on $n$ Boolean variables, and write $\mathrm{coeffs}(T) :
\mathcal{P}(\{0,\dots,n-1\}) \to \bbr$ for its recursively defined coefficient function.
If $k$ is a natural number with $\mathrm{depth}(T) \le k$, then for every set $S$ of
variables the scaled coefficient $\mathrm{coeffs}(T)(S)\cdot 2^{k}$ is an integer; that
is, there exists $m \in \bbz$ with $\mathrm{coeffs}(T)(S)\cdot 2^{k} = m$.
\end{lemma}

\begin{lemma}[Dyadic granularity of a decision tree's coefficients]
\lean{DecisionTree.coeffs_granular}
\statementsource{boolean-ch03-spectral-learning-restrictions}{pdf-20e15a3d3d5e-p006-b003}
\label{DecisionTree.coeffs_granular}
\leanok
\uses{DecisionTree, DecisionTree.coeffs, DecisionTree.coeffs_mul_two_pow_int, DecisionTree.depth}
\difficulty{2}
Let $T$ be a decision tree on $n$ Boolean variables, and let $S \subseteq
\{0,\dots,n-1\}$ be a frequency set. Then the coefficient $\mathrm{coeffs}\,T\,S$ is a
dyadic rational with denominator $2^{\mathrm{depth}(T)}$: there exists an integer $m \in
\bbz$ such that \[ \mathrm{coeffs}\,T\,S = \frac{m}{2^{\mathrm{depth}(T)}}. \]
Equivalently, every such coefficient is an integer multiple of $2^{-\mathrm{depth}(T)}$.
\end{lemma}

\begin{lemma}[Size bounded by two to the depth]
\lean{DecisionTree.size_le_two_pow_depth}
\label{DecisionTree.size_le_two_pow_depth}
\leanok
\uses{DecisionTree, DecisionTree.depth, DecisionTree.size}
\difficulty{4}
Let $T$ be a decision tree on $n$ Boolean variables, with size (number of leaves)
$\mathrm{size}(T)$ and depth $\mathrm{depth}(T)$. Then $\mathrm{size}(T) \le
2^{\mathrm{depth}(T)}$; that is, a decision tree of depth $k$ has at most $2^{k}$
leaves.
\end{lemma}

\begin{lemma}[Recovery of coefficients from a character expansion]
\lean{DecisionTree.fourierCoeff_sum_chiS}
\label{DecisionTree.fourierCoeff_sum_chiS}
\leanok
\uses{BooleanAnalysis.BoolCube, BooleanAnalysis.chiS, BooleanAnalysis.fourierCoeff, BooleanAnalysis.fourier_coeff_chi, BooleanAnalysis.innerProduct, BooleanAnalysis.uniformWeight}
\difficulty{4}
Fix $n$, let $c$ assign a real number $c(S)$ to each subset $S \subseteq [n]$, and
consider the Boolean function $f(x) = \sum_{S \subseteq [n]} c(S)\,\chi_S(x)$, where
$\chi_S$ is the Walsh--Fourier character of $S$. Then for every $T \subseteq [n]$ the
Fourier--Walsh coefficient of $f$ at frequency $T$ recovers the corresponding
coefficient, namely $\hat f(T) = c(T)$.
\end{lemma}

\begin{theorem}[Fourier coefficients of a decision tree]
\lean{DecisionTree.fourierCoeff_signEval}
\label{DecisionTree.fourierCoeff_signEval}
\leanok
\uses{BooleanAnalysis.chiS, BooleanAnalysis.fourierCoeff, DecisionTree, DecisionTree.coeffs, DecisionTree.fourierCoeff_sum_chiS, DecisionTree.signEval, DecisionTree.signEval_eq_sum_coeffs}
\difficulty{2}
Let $T$ be a decision tree on $n$ Boolean variables, and let $\sigma\!\circ\! T$ denote
its sign-encoded function $x \mapsto \sigma(T(x))$ on the Boolean hypercube $\{0,1\}^n$.
Then for every frequency set $S \subseteq [n]$, the Fourier–Walsh coefficient of this
function at $S$ agrees with the recursively defined tree coefficient:
\[
  \widehat{\sigma\!\circ\! T}(S) \;=\; \mathrm{coeffs}(T)(S).
\]
\end{theorem}

\begin{theorem}[Fourier degree of a decision tree is at most its depth]
\lean{DecisionTree.degree_le_depth}
\label{DecisionTree.degree_le_depth}
\leanok
\statementsource{boolean-ch03-spectral-learning-restrictions}{pdf-20e15a3d3d5e-p006-b003}
\uses{BooleanAnalysis.has_degree_at_most, DecisionTree, DecisionTree.coeffs_eq_zero_of_depth_lt, DecisionTree.depth, DecisionTree.fourierCoeff_signEval, DecisionTree.signEval}
\difficulty{3}
Let $T$ be a decision tree on $n$ Boolean variables, and let $f\colon\{0,1\}^n\to\bbr$
be the $\pm 1$-valued function it computes, so that $f(x)=\sigma(T(x))$ where $\sigma$
sends $\mathrm{false}$ to $1$ and $\mathrm{true}$ to $-1$. Then $f$ has Fourier degree
at most the depth of $T$: for every $S\subseteq\{1,\dots,n\}$ with $\hat f(S)\ne 0$, one
has $\abs{S}\le\operatorname{depth}(T)$.
\end{theorem}

\begin{theorem}[Spectral one-norm bounded by decision tree size]
\lean{DecisionTree.spectral_one_norm_le}
\label{DecisionTree.spectral_one_norm_le}
\leanok
\statementsource{boolean-ch03-spectral-learning-restrictions}{pdf-20e15a3d3d5e-p006-b003}
\uses{BooleanAnalysis.fourierCoeff, DecisionTree, DecisionTree.coeffs, DecisionTree.fourierCoeff_signEval, DecisionTree.signEval, DecisionTree.size, DecisionTree.sum_abs_coeffs_le}
\difficulty{2}
Let $T$ be a decision tree on $n$ Boolean variables, and let $g : \{0,1\}^n \to
\{-1,1\}$ be the sign-encoded function it computes, sending an input $x$ to
$\sigma(T(x))$ (with $\mathrm{false}\mapsto 1$ and $\mathrm{true}\mapsto -1$). Then the
sum of the absolute values of its Fourier--Walsh coefficients is bounded by the size of
$T$, i.e. its number of leaves:
\[
  \sum_{S \subseteq [n]} \bigl|\hat g(S)\bigr| \;\le\; \mathrm{size}(T),
\]
where the sum ranges over all subsets $S$ of $\{1,\dots,n\}$.
\end{theorem}

\begin{theorem}[Granularity of decision-tree Fourier coefficients]
\lean{DecisionTree.fourierCoeff_granular}
\label{DecisionTree.fourierCoeff_granular}
\leanok
\statementsource{boolean-ch03-spectral-learning-restrictions}{pdf-20e15a3d3d5e-p006-b003}
\uses{BooleanAnalysis.fourierCoeff, DecisionTree, DecisionTree.coeffs_granular, DecisionTree.depth, DecisionTree.fourierCoeff_signEval, DecisionTree.signEval}
\difficulty{1}
Let $T$ be a decision tree on $n$ Boolean variables, let $f\colon\{0,1\}^n\to\{-1,1\}$
be the $\pm1$-valued function it computes under the sign encoding $\mathrm{false}\mapsto
1$, $\mathrm{true}\mapsto -1$, and let $d$ be the depth of $T$. Then for every set of
coordinates $S\subseteq[n]$ there is an integer $m$ such that the Fourier--Walsh
coefficient of $f$ at frequency $S$ satisfies
\[
  \hat f(S) \;=\; \frac{m}{2^{d}}.
\]
That is, every Fourier coefficient of $f$ is an integer multiple of $2^{-d}$.
\end{theorem}

\begin{theorem}[Fourier sparsity bound for decision trees]
\lean{DecisionTree.sparsity_le}
\label{DecisionTree.sparsity_le}
\leanok
\statementsource{boolean-ch03-spectral-learning-restrictions}{pdf-20e15a3d3d5e-p006-b003}
\uses{BooleanAnalysis.fourierCoeff, DecisionTree, DecisionTree.depth, DecisionTree.fourierCoeff_granular, DecisionTree.signEval, DecisionTree.size, DecisionTree.spectral_one_norm_le}
\difficulty{5}
Let $T$ be a decision tree on $n$ Boolean variables, and let $g \colon \{0,1\}^n \to
\{-1,1\}$ be the sign-encoded function it computes. Then the number of frequencies $S
\subseteq [n]$ at which the Fourier--Walsh coefficient of $g$ is nonzero satisfies
\[
\abs{\{\, S \subseteq [n] : \hat g(S) \neq 0 \,\}} \;\le\; \mathrm{size}(T)\cdot
2^{\,\mathrm{depth}(T)},
\]
where $\mathrm{size}(T)$ is the number of leaves of $T$ and $\mathrm{depth}(T)$ is its
maximum root-to-leaf path length.
\end{theorem}

\begin{theorem}[Fourier sparsity bound for decision trees]
\lean{DecisionTree.sparsity_le_four_pow}
\label{DecisionTree.sparsity_le_four_pow}
\leanok
\statementsource{boolean-ch03-spectral-learning-restrictions}{pdf-20e15a3d3d5e-p006-b003}
\uses{BooleanAnalysis.fourierCoeff, DecisionTree, DecisionTree.depth, DecisionTree.signEval, DecisionTree.size, DecisionTree.size_le_two_pow_depth, DecisionTree.sparsity_le}
\difficulty{2}
Let $T$ be a decision tree on $n$ Boolean variables, and let
$f\colon\{0,1\}^n\to\{-1,1\}$ be its sign-encoded tree function. Then the number of
subsets $S\subseteq[n]$ whose Fourier--Walsh coefficient $\hat f(S)$ is nonzero is at
most $4^{\,\mathrm{depth}(T)}$.
\end{theorem}

\begin{lemma}[Attainment of the decision-tree depth]
\lean{DecisionTree.exists_dtree_of_dtDepth}
\label{DecisionTree.exists_dtree_of_dtDepth}
\leanok
\statementsource{boolean-ch03-spectral-learning-restrictions}{pdf-20e15a3d3d5e-p006-b001}
\uses{DecisionTree, DecisionTree.depth, DecisionTree.eval, dtDepth}
\difficulty{1}
For every Boolean function $f : (\mathrm{Fin}\,n \to \mathrm{Bool}) \to \mathrm{Bool}$,
there is a decision tree $T$ on the $n$ Boolean variables whose depth is at most the
decision-tree depth of $f$ and which computes $f$, meaning $T$ evaluates to $f(x)$ on
every input $x$.
\end{lemma}

\begin{theorem}[Degree bounded by decision-tree depth]
\lean{DecisionTree.degree_le_dtDepth}
\label{DecisionTree.degree_le_dtDepth}
\leanok
\statementsource{boolean-ch03-spectral-learning-restrictions}{pdf-20e15a3d3d5e-p006-b003}
\uses{BooleanAnalysis.boolToSign, BooleanAnalysis.has_degree_at_most, DecisionTree.degree_le_depth, DecisionTree.exists_dtree_of_dtDepth, DecisionTree.signEval, dtDepth}
\difficulty{3}
Let $f : \{0,1\}^n \to \{0,1\}$ be a Boolean-valued function, and let $D(f)$ denote its
decision-tree depth, the least $d$ for which some decision tree of depth at most $d$
computes $f$. Let $g : \{0,1\}^n \to \bbr$ be its sign encoding, $g(x) = \sigma(f(x))$,
where $\sigma$ sends $0$ to $1$ and $1$ to $-1$. Then $g$ has degree at most $D(f)$:
every set $S$ with $\hat g(S) \neq 0$ satisfies $\abs{S} \le D(f)$.
\end{theorem}
